% !TEX program = xelatex
% CUMCM 论文骨架（按 2025 国赛 A 题实战反馈重构；2026-08 增补图表插入时机 / 双引擎 / 专业彩色配色）
% 版式：A4 白纸、四边页边距 >= 2.5cm、页码页脚中部、摘要起从 1 连续编号、图表居中、图注中文
% 摘要页与正文统一排版体系：标题（三号黑体居中）→「摘要」（四号黑体居中，与正文一级标题同字号）
%   → 正文一级标题「一、问题重述」黑体四号 + 中文数字，二级/三级保留阿拉伯编号（1.1 / 1.1.1）
% 图表规范（详见 references/04-图表规范.md）：
%   - 每个已求解子问题默认 >= 3 张图：工作机理示意图 + 算法流程图 + 核心结果图；验证/灵敏度另加
%   - 图前有引导文字、图后有解释文字；连续 2-3 张图表间必须有文字
%   - 配色：专业彩色（色盲友好色板）或灰度均可，灰度可区分；禁止彩虹渐变
%   - 引擎：求解与出图可用 Python matplotlib 或 MATLAB（exportgraphics），规范见 references/04、10
\documentclass[11pt]{ctexart}
\usepackage[a4paper,top=2.5cm,bottom=2.5cm,left=2.5cm,right=2.5cm]{geometry}
\usepackage{amsmath,amssymb}
\usepackage{graphicx}
\usepackage{booktabs}
\usepackage{caption}
\usepackage{float}
\usepackage{listings}% 附录 B 代码展示（每问求解/模型代码 + 逐行行号，不含生图代码）
\usepackage{fancyhdr}
\usepackage[hidelinks]{hyperref}
\usepackage{enumitem}
\captionsetup{font=small,labelfont=bf,justification=centering}

% ---- 页码：页脚中部，阿拉伯数字从摘要页 1 连续编号 ----
\pagestyle{fancy}
\fancyhf{}
\fancyfoot[C]{\thepage}
\renewcommand{\headrulewidth}{0pt}

% ---- 章节样式统一（与摘要页同一排版体系）----
% 一级标题：黑体四号 + 中文数字「一、问题重述」（aftername 置空去掉编号后多余间距）
% 二级/三级：黑体小四，保留阿拉伯编号（1.1 / 1.1.1）
\ctexset{
  section = {
    format = \centering\heiti\zihao{4},
    name   = {,、},
    aftername = {},
    beforeskip = 1.2ex,
    afterskip  = 1.0ex
  },
  subsection   = {format = \heiti\zihao{-4}},
  subsubsection= {format = \heiti\zihao{-4}}
}
\renewcommand\thesection{\chinese{section}}
\renewcommand\thesubsection{\arabic{section}.\arabic{subsection}}
\renewcommand\thesubsubsection{\arabic{section}.\arabic{subsection}.\arabic{subsubsection}}

\begin{document}

% ========== 摘要（单独一页，<= 1 页）==========
% 结构：背景段 + 每个问题各一段 + 关键词行（无英文翻译）
% 每问段落七要素：①情形及与前问关系 ②建模（决策变量+目标函数）③算法点名+亮点
%                ④最优策略数值（6~8 位小数）⑤结果 ⑥交叉验证/相对误差 ⑦附件说明
% 标题块：与「摘要」标题统一用「居中 + 黑体」体系（标题三号、摘要四号），
% 不再用 \maketitle（article 样式 + 空 author/date 留白，与正文 \section 割裂）
\begin{center}
{\zihao{3}\heiti\bfseries <赛题名称>}
\end{center}

\vspace{1.0em}
\begin{center}
{\zihao{4}\heiti 摘要}
\end{center}

\vspace{0.6em}
<背景段：本文针对……赛题，建立了……模型。问题难点在于……>

<问题一：针对问题一，建立……模型，以……为决策变量、……为目标函数，采用……算法求解，得到……（6~8 位小数）。为验证可靠性，用……交叉验证，相对误差为……。其余数据见附件 result1.xlsx。>

<问题二：……（同上七要素）>

<问题三：……>

<问题四：……>

<问题五：……>

\vspace{0.6em}
\noindent\textbf{关键词：}关键词1；关键词2；关键词3；关键词4
\setcounter{page}{1}
\newpage

% ============================================================
% 图表插入模板（复制到目标位置后取消注释并替换占位内容）
%
% 图：图前引导文字 -> \begin{figure} -> 图后解释文字
% 图注要点：①展示什么 ②关键趋势/结果 ③验证含义(若是验证图)；禁止「支撑哪问/数据来源」流程解释
% 题注排版：单独一行并居中；caption 内【不写】「图 N：」/「表 N：」前缀——编号由 \caption 自动生成，
%            手写前缀会造成「图 10: 图 10：…」图号重复
% \begin{figure}[H]
%   \centering
%   \includegraphics[width=0.85\textwidth]{../figures/fig_q1_result.png}  % LaTeX 与 Word 均用 PNG
%   \caption{<展示什么>，结果表明<关键趋势/结果>。}
%   \label{fig:q1_result}
% \end{figure}
%
% 三线表（结果表/可行性表/灵敏度表/评价表）：
% \begin{table}[H]
%   \centering
%   \caption{<表注>}
%   \label{tab:q1}
%   \begin{tabular}{ccc}
%     \toprule 列1 & 列2 & 列3 \\
%     \midrule
%     <值> & <值> & <值> \\
%     \bottomrule
%   \end{tabular}
% \end{table}
% ============================================================

% ========== 一、问题重述 ==========
\section{问题重述}
\subsection{问题背景}
<背景：问题涉及的物理原理/工程背景 + 军事/实际价值 + 本文任务总述。>

\subsection{问题提出}
<先总述已知条件，再按本质归类子问题（如：问题一、二为单机单弹求值→优化；问题三为多弹时序规划；问题四为多机空间协同；问题五为任务分配+资源调度）。>
% ⚠️ 禁止写「结果填入 resultX.xlsx」这类话术，全部用建模化措辞。

% 问题列表：直接写「问题一：」…文本编号，不用 enumerate 自动编号（避免「1. 问题 1」重复编号）
\begin{itemize}[leftmargin=2em,itemsep=1pt,label={}]
\item \textbf{问题一}：<……>
\item \textbf{问题二}：<……>
\item \textbf{问题三}：<……>
\item \textbf{问题四}：<……>
\item \textbf{问题五}：<……>
\end{itemize}

% ========== 二、问题分析 ==========
\section{问题分析}
% 结构（2024A 复盘升级，范式详见 references/05-论文写作.md）：2.1 总体分析 + 2.2~2.6 逐问五步法；篇幅 1.5~2 页
<总起段：一句话总述五问共用内核与依赖关系，引出技术路线图。>

% 🔎 节首放 1 张全题技术路线图（fig_route_overview.png，逐问解题路径总览），图前有引导文字：
% \begin{figure}[H]
%   \centering
%   \includegraphics[width=1.0\textwidth]{fig_00_route_overview.png}
%   \caption{<全题技术路线图：五问的建模目标、求解路线、输出与难点突破，标注共享内核与依赖关系。>}
%   \label{fig:route}
% \end{figure}

% ---- 2.1 总体分析（新增分节，必写）----
% 三段式：①本质归类（一句话系统类型 + 落点 + 为什么不能降维）
%         ②核心矛盾/耦合主线（2~4 条，每条 = 一句话机理 + 对应建模手段，主线数 = 后文模型模块数）
%         ③量级观察与风险锁定（用题给数据估算量级 → 预测关键现象，后文结果处呼应）
%         （可选）④分层结构与递进关系（子问题信息流与复用）
\subsection{总体分析}
<本质归类：本题在本质上是一个"<系统>"的<类型>问题：<对象与约束>……；落点不是<单点/局部>描述，而是<系统级统一求解>。与<更简单模型>不同，须同时处理 N 类本质耦合：其一是<机理>→ 对应<建模手段>；其二是……；其三是……。>
<量级观察：由题给数据估算……（总长/跨越圈数/间隙与尺寸对比/时间尺度），可预测<关键现象>，因此<后文某问>正是对这一过程的刻画。>

% ---- 2.2~2.6 逐问分节：五步法 ----
% 每节五步：①与前问关系 ②分步求解路线 ③机理层"为什么"（本质原因，非步骤复述）
%           ④本问特有难点 ⑤落点（建什么模型）；（推荐）⑥数值佐证（题给参数或已算结果）
\subsection{问题一分析}
<① 与前问关系；② 求解路线；③ 机理为什么……；④ 难点；⑤ 落点。数值佐证：……。>

\subsection{问题二分析}
<同上五步。>

\subsection{问题三分析}
<同上五步。>

\subsection{问题四分析}
<同上五步。>

\subsection{问题五分析}
<同上五步。>

% ========== 三、模型假设 ==========
\section{模型假设}
\begin{enumerate}[leftmargin=2em,itemsep=1pt]
\item <假设 1（含取值依据与灵敏度考察说明）>
\item <假设 2>
\item <假设 3>
\end{enumerate}

% ========== 四、符号说明 ==========
\section{符号说明}
% 规则：删除表名/表号/label；每个符号单独一行，禁止同行两个符号
\begin{center}
\begin{tabular}{lll}
\toprule
符号 & 含义 & 单位 \\
\midrule
$t$ & <含义> & <单位> \\
$x$ & <含义> & <单位> \\
$y$ & <含义> & <单位> \\
$v$ & <含义> & <单位> \\
$\theta$ & <含义> & <单位> \\
\bottomrule
\end{tabular}
\end{center}

% ========== 五、模型准备 ==========
\section{模型准备}
\subsection{坐标系与变量约定}
<建立坐标系，约定符号与单位，说明维度/方向约定。>

\subsection{态势可视化与观察分析}
<画态势图（俯视图/侧视图），从图中观察几何关系，得出建模动机。>
% 🔎 插入时机：此处放「态势/坐标系示意图」（fig_problem_overview）+ 全题工作机理示意图（fig_mechanism_overview，可并入态势图）——先画图再建模
% 图前引导：<说明为什么要先看图、图中观察到什么几何关系>

\subsection{运动轨迹方程}
<推导基础运动方程：如抛体运动、匀速直线运动、云团下沉等。>

\subsection{关键判定条件}
<推导核心判定式：如视线遮挡判定、有效范围判定、约束条件。>

% ========== 六、模型的建立与求解 ==========
\section{模型的建立与求解}
% 每问设副标题（不单独大标题）；建模四段式 + 求解三段式；图是「观察→动机→算法」的起点
\subsection{问题一：<名称>}
\subsubsection{模型的建立}
<四段式：①情形及与前问关系 ②决策变量与目标函数 ③约束条件 ④建立公式/方程。>
% 优化模型展示：目标函数 + s.t. + 大括号约束组（思考/推导过程逐步写在正文，禁止只贴最终模型）
% \begin{align}
%   \max_{x}\quad & f(x) \\
%   \text{s.t.}\quad & \left\{
%   \begin{array}{ll}
%     g_i(x)\le 0, & i=1,2,\dots,m,\\[2pt]
%     h_j(x)=0, & j=1,2,\dots,n,\\[2pt]
%     x\in \mathcal{X}.
%   \end{array}
%   \right.
% \end{align}
\begin{equation}
\label{eq:q1}
<公式>
\end{equation}

\subsubsection{模型的求解}
<三段式：①算法（点名+亮点） ②数值求解 ③结果与分析。>
<结果数值保留 6~8 位小数。>
% 🔎 插入时机：每问默认 3 张图——工作机理示意图（fig_q1_mechanism）+ 算法流程图（fig_q1_flow）+ 核心结果图（fig_q1_result）；
%              有验证/误差/约束/灵敏度时再加 1 张检查图（fig_q1_validation）
% 图前引导：<为什么需要这张图、它检验或回答什么问题>；图后解释：<关键数值/趋势/对比，回到题目要求>

\subsection{问题二：<名称>}
\subsubsection{模型的建立}
<同上四段式。>
\subsubsection{模型的求解}
<同上三段式。>
% 🔎 插入 fig_q2_mechanism（工作机理示意图）+ fig_q2_flow（算法流程图）+ fig_q2_result（结果图）（+ 验证/灵敏度图）

\subsection{问题三：<名称>}
\subsubsection{模型的建立}
\subsubsection{模型的求解}
% 🔎 插入 fig_q3_mechanism（工作机理示意图）+ fig_q3_flow（算法流程图）+ fig_q3_result（结果图）（+ 验证/灵敏度图）

\subsection{问题四：<名称>}
\subsubsection{模型的建立}
\subsubsection{模型的求解}
% 🔎 插入 fig_q4_mechanism（工作机理示意图）+ fig_q4_flow（算法流程图）+ fig_q4_result（结果图）（+ 验证/灵敏度图）

\subsection{问题五：<名称>}
\subsubsection{模型的建立}
\subsubsection{模型的求解}
% 🔎 插入 fig_q5_mechanism（工作机理示意图）+ fig_q5_flow（算法流程图）+ fig_q5_result（结果图）（+ 验证/灵敏度图）

% ========== 七、模型的检验 ==========
\section{模型的检验}
\subsection{灵敏度分析}
<关键参数 ±10%~±50% 扰动，结果变化范围与结论。>
% 🔎 插入时机：独立章节放「参数扫描曲线/热图/Tornado 图」（fig_sensitivity），不塞进结果段
% 图后解释：<哪个参数最敏感、变化方向、对结论的影响>

\subsection{稳定性与误差分析}
<网格/步长收敛性、与解析解或实测对比、相对误差。>
% 🔎 可插入收敛性图（fig_convergence）或 解析 vs 数值 同轴对比图（fig_compare）

% ========== 八、模型的评价与推广 ==========
\section{模型的评价与推广}
% 说明：本节用文字分节陈述（优点/缺点/推广），模板不默认插入表格；
%       仅当确有「多方案/多维对比」价值时才可另加三线表（表注居中，编号由 LaTeX 自动生成）
\subsection{模型的优点}
% 规则：优点条数 > 缺点条数
\begin{enumerate}[leftmargin=2em,itemsep=1pt]
\item <优点 1>
\item <优点 2>
\item <优点 3>
\end{enumerate}
\subsection{模型的缺点}
% 规则：缺点用「代价型」措辞（如「以……计算代价换取……」「在……条件下适用」）
\begin{enumerate}[leftmargin=2em,itemsep=1pt]
\item <以……计算代价换取……>
\item <在……假设下适用，偏离时需修正>
\end{enumerate}
\subsection{模型的推广}
<推广到其他场景/问题。>

% ========== AI 使用声明（参考文献之前）==========
% 独立成段、与正文字号一致；\cite 所用 AI 模型，参考文献中同步加该模型版本条目（见 references/05）
% \noindent\textbf{声明：}本参赛队在竞赛过程中使用了 AI 工具（<模型名> 大语言模型\cite{<ai-key>}），主要用于语言润色、代码调试和排版等，详细使用情况见支撑材料「AI使用情况」。

% ========== 参考文献 ==========
% 规则：只保留 thebibliography；不少于 8 条学术文献；每条至少被正文 \cite 一次；删除官网/组委会条目
% 参考文献标题：居中（\renewcommand{\refname}{} 抑制自动标题后手动居中输出）
\renewcommand{\refname}{}
\begin{center}
{\zihao{4}\heiti 参考文献}
\end{center}
\vspace{1.0ex}
\begin{thebibliography}{99}
\bibitem{ref1} <作者. 标题. 期刊, 年份, 卷(期): 页码.>
\bibitem{ref2} <作者. 标题. 会议/期刊, 年份.>
\bibitem{ref3} <作者. 标题. 期刊, 年份.>
\bibitem{ref4} <作者. 标题. 期刊, 年份.>
\bibitem{ref5} <作者. 标题. 期刊, 年份.>
\bibitem{ref6} <作者. 标题. 期刊, 年份.>
\bibitem{ref7} <作者. 标题. 期刊, 年份.>
\bibitem{ref8} <作者. 标题. 期刊, 年份.>
\end{thebibliography}
% 正文必须对每条文献至少 \cite{refN} 一次（用 compliance_check.py 自动检查）

% ========== 附录（不编序号）==========
\begin{center}
{\Large\bfseries 附录}
\end{center}

\section*{附录 A　支撑材料目录}
<按 collect\_materials.py 生成的支撑材料清单粘贴：源程序、中间图表、结果数据。>

\section*{附录 B　代码（Python / MATLAB）}
% 要求：展示支撑材料（materials/源程序/）中的【每问求解代码 + 模型代码】，不含生图代码；
%       画图脚本统一放 code/plots/ 或命名 plot_*.py / fig_*.py / draw_*.m（collect_materials.py 自动排除）；
%       每个代码文件一个小节：先写「代码名」，再用 \lstinputlisting 逐行标明行号列出完整代码。
% 用法：1) 导言区已加载 listings 包（\lstset 配置见下）；
%       2) 求解/模型代码有多少个文件就列多少个小节，全部展示，不许只列文件名；
%       3) 删除下面示例前的 % 并按实际文件路径替换。
\lstset{
  basicstyle=\ttfamily\small,
  numbers=left,
  numberstyle=\tiny,
  stepnumber=1,
  breaklines=true,
  showstringspaces=false,
  frame=single,
  language=Python
}

% 示例：每个源码文件一个小节（代码名 + 逐行行号）。取消注释并替换路径后启用：
% \subsection*{B.1　code/main.py}
% \lstinputlisting{../code/main.py}
%
% \subsection*{B.2　code/utils.py}
% \lstinputlisting{../code/utils.py}
%
% MATLAB 同理：\subsection*{B.3　code/opt.m} \lstinputlisting{../code/opt.m}

\end{document}
